%%%%%%%%%%%%%%%%%
% This is an sample CV template created using altacv.cls
% (v1.7.4, 30 July 2025) written by LianTze Lim (liantze@gmail.com). Compiles with pdfLaTeX, XeLaTeX and LuaLaTeX.
%
%% It may be distributed and/or modified under the
%% conditions of the LaTeX Project Public License, either version 1.3
%% of this license or (at your option) any later version.
%% The latest version of this license is in
%%    http://www.latex-project.org/lppl.txt
%% and version 1.3 or later is part of all distributions of LaTeX
%% version 2003/12/01 or later.
%%%%%%%%%%%%%%%%

%% Use the "normalphoto" option if you want a normal photo instead of cropped to a circle
% \documentclass[10pt,a4paper,withhyper,normalphoto]{altacv}

\documentclass[10pt,a4paper,withhyper]{../altacv}
%% AltaCV uses the fontawesome5 and simpleicons packages.
%% See http://texdoc.net/pkg/fontawesome5 and http://texdoc.net/pkg/simpleicons for full list of symbols.

% Change the page layout if you need to
\geometry{left=1.25cm,right=1.25cm,top=1.5cm,bottom=1.5cm,columnsep=1.2cm}

% The paracol package lets you typeset columns of text in parallel
\usepackage{paracol}

% Change the font if you want to, depending on whether
% you're using pdflatex or xelatex/lualatex
% WHEN COMPILING WITH XELATEX PLEASE USE
% xelatex -shell-escape -output-driver="xdvipdfmx -z 0" sample.tex
\iftutex
  % If using xelatex or lualatex:
  \setmainfont{Roboto Slab}
  \setsansfont{Lato}
  \renewcommand{\familydefault}{\sfdefault}
\else
  % If using pdflatex:
  \usepackage[rm]{roboto}
  \usepackage[defaultsans]{lato}
  % \usepackage{sourcesanspro}
  \renewcommand{\familydefault}{\sfdefault}
\fi

% Change the colours if you want to
\definecolor{SlateGrey}{HTML}{2E2E2E}
\definecolor{LightGrey}{HTML}{666666}
\definecolor{DarkPastelRed}{HTML}{450808}
\definecolor{PastelRed}{HTML}{8F0D0D}
\definecolor{GoldenEarth}{HTML}{E7D192}
\colorlet{name}{black}
\colorlet{tagline}{PastelRed}
\colorlet{heading}{DarkPastelRed}
\colorlet{headingrule}{GoldenEarth}
\colorlet{subheading}{PastelRed}
\colorlet{accent}{PastelRed}
\colorlet{emphasis}{SlateGrey}
\colorlet{body}{LightGrey}

% Change some fonts, if necessary
\renewcommand{\namefont}{\Huge\rmfamily\bfseries}
\renewcommand{\personalinfofont}{\footnotesize}
\renewcommand{\cvsectionfont}{\LARGE\rmfamily\bfseries}
\renewcommand{\cvsubsectionfont}{\large\bfseries}


% Change the bullets for itemize and rating marker
% for \cvskill if you want to
\renewcommand{\cvItemMarker}{{\small\textbullet}}
\renewcommand{\cvRatingMarker}{\faCircle}
% ...and the markers for the date/location for \cvevent
% \renewcommand{\cvDateMarker}{\faCalendar*[regular]}
% \renewcommand{\cvLocationMarker}{\faMapMarker*}


% If your CV/résumé is in a language other than English,
% then you probably want to change these so that when you
% copy-paste from the PDF or run pdftotext, the location
% and date marker icons for \cvevent will paste as correct
% translations. For example Spanish:
% \renewcommand{\locationname}{Ubicación}
% \renewcommand{\datename}{Fecha}


%% Use (and optionally edit if necessary) this .tex if you
%% want to use an author-year reference style like APA(6)
%% for your publication list
% \input{pubs-authoryear}

%% Use (and optionally edit if necessary) this .tex if you
%% want an originally numerical reference style like IEEE
%% for your publication list
\usepackage[backend=biber,style=ieee,sorting=ydnt,defernumbers=true]{biblatex}
%% For removing numbering entirely when using a numeric style
\setlength{\bibhang}{1.25em}
\DeclareFieldFormat{labelnumberwidth}{\makebox[\bibhang][l]{\itemmarker}}
\setlength{\biblabelsep}{0pt}
\defbibheading{pubtype}{\cvsubsection{#1}}
\renewcommand{\bibsetup}{\vspace*{-\baselineskip}}
\AtEveryBibitem{%
  \iffieldundef{doi}{}{\clearfield{url}}%
}


%% sample.bib contains your publications
\addbibresource{sample.bib}

\begin{document}
\name{Guanlin Wang}
\tagline{Computer Vision / AI Engineer}
%% You can add multiple photos on the left or right
\photoR{2.8cm}{../../images/gl_photo}
% \photoL{2.5cm}{Yacht_High,Suitcase_High}

\personalinfo{%
  % Not all of these are required!
  \email{g.wang@mail.de}
  \phone{+49 151 56312173}
  \mailaddress{Alfred-Messel-Weg 38, 64287, Darmstadt}
  \printinfo{\faLinkedin}{Guanlin Wang}[https://www.linkedin.com/in/guanlin-wang-737354206]
}

\makecvheader
%% Depending on your tastes, you may want to make fonts of itemize environments slightly smaller
% \AtBeginEnvironment{itemize}{\small}

%% Set the left/right column width ratio to 6:4.
\columnratio{0.3}

% Start a 2-column paracol. Both the left and right columns will automatically
% break across pages if things get too long.
\begin{paracol}{2}

  \cvsection{FÄHIGKEITEN}

  % Don't overuse these \cvtag boxes — they're just eye-candies and not essential. If something doesn't fit on a single line, it probably works better as part of an itemized list (probably inlined itemized list), or just as a comma-separated list of strengths.
  \cvsubsection{Computer Vision}
  
  \vspace{0.5em}

  \cvtag{Computer Vision}

  \vspace{0.5em}

  \cvtag{Bildverarbeitung}

  \vspace{0.5em}

  \cvtag{Deep Learning}

  \divider\smallskip

  \cvsubsection{Computer Vision Stack}

  \vspace{0.5em}

  \cvtag{OpenCV}

  \vspace{0.5em}

  \cvtag{PyTorch}
  
  \vspace{0.5em}

  \cvtag{ONNX}

  \vspace{0.5em}

  \cvtag{YOLO}

  \divider\smallskip

  \cvsubsection{Software Development}

  \vspace{0.5em}

  \cvtag{Python}
  
  \vspace{0.5em}

  \cvtag{C++}
  
  \vspace{0.5em}

  \cvtag{C\#}

  \vspace{0.5em}

  \cvtag{REST-API}
  
  \vspace{0.5em}
  
  \cvtag{FastAPI}

  \vspace{0.5em}

  \cvtag{ASP.NET Core}

  \vspace{0.5em}

  \cvtag{.NET Blazor Server}

  \divider\smallskip

  \cvsubsection{Tools}

  \vspace{0.5em}

  \cvtag{Docker}

  \vspace{0.5em}

  \cvtag{Git}
  
  \vspace{0.5em}
  
  \cvtag{Linux}


  \newpage
  \cvsection{SPRACHEN}

  \vspace{0.5em}

  \cvtag{Chinesisch (Muttersprache)}

  \vspace{0.5em}

  \cvtag{Deutsch (C1)}

  \vspace{0.5em}

  \cvtag{Englisch (B2)}


  %% Yeah I didn't spend too much time making all the
  %% spacing consistent... sorry. Use \smallskip, \medskip,
  %% \bigskip, \vspace etc to make adjustments.
  \medskip

  % use ONLY \newpage if you want to force a page break for
  % ONLY the current column
  \newpage

  %% Switch to the right column. This will now automatically move to the second
  %% page if the content is too long.
  \switchcolumn
  \cvsection{BERUFSERFAHRUNG}

  \cvevent{Computer Vision / AI Entwickler}{Innomatik AG}{06/2021 -- Heute}{Bensheim}

  \vspace{1em}

  % ocr
  \cvsubsubsection{OCR-Erkennungspipeline für reale Bilder und technische Zeichnungen}

  \vspace{0.5em}

  \begin{itemize}
    \item Entwicklung einer modularen OCR-Pipeline für reale Bilder und technische Zeichnungen mit \textbf{Text Detection} und \textbf{Text Recognition}
    \item Aufbau einer vollständigen Trainings- und Verarbeitungspipeline einschließlich Datenvorbereitung, Annotation, Modelltraining sowie Bildvor- und Nachverarbeitung
    \item Deployment einer lokalen Desktop-Anwendung sowie als \textbf{REST-API} mit \textbf{FastAPI} und \textbf{Docker}
  \end{itemize}

  \vspace{0.5em}

  \textbf{Technologien}: Python, OpenCV, PyTorch, PaddleOCR, FastAPI, REST-API, Docker, uv, Git

  \vspace{0.5em}

  % object detection
  \cvsubsubsection{Object-Detection-System für reale Bildszenarien}

  \vspace{0.5em}

  \begin{itemize}
    \item Entwicklung eines \textbf{YOLO}-basierten \textbf{Object-Detection}-Systems mit automatisierter Trainingspipeline und \textbf{Multi-Scale-Inferenz}
    \item Aufbau einer automatisierten Trainingspipeline einschließlich Datensatzaufbereitung, Annotation, eigener Datenaugmentation, Modelltraining und Modellevaluation für ein \textbf{YOLOv12-L}-Modell (\textbf{21 Klassen, ~1.500 Bilder, mAP@50-95: 85-90 \%})
    \item Deployment der \textbf{ONNX}-Modelle als \textbf{C\# ASP.NET Core WebAPI} mit \textbf{Docker} sowie Integration in eine \textbf{.NET Blazor-Server-Anwendung}
  \end{itemize}

  \vspace{0.5em}

  \textbf{Technologien}: Python, YOLOv12, ONNX, C\#, ASP.NET Core WebAPI, Blazor Server, Docker, uv, Git

  \vspace{0.5em}

  % visual location
  \cvsubsubsection{Visuelle Lokalisierung — Kamerarelokalisierung für Digital-Twin-AR}

  \vspace{0.5em}

  \begin{itemize}
    \item Entwicklung einer \textbf{Visual-Localization}-Pipeline zur \textbf{Kamerapose-Schätzung} auf Basis von \textbf{Feature Matching} und \textbf{2D-3D-Korrespondenzen}
    \item Entwicklung einer robusten \textbf{Multi-Object-/Multi-View-Pipeline} zur Erhöhung der Lokalisierungsrobustheit mit \textbf{klassischen Feature-Matching}-Verfahren
    \item Evaluierung und Integration moderner \textbf{Deep-Learning}-Verfahren (\textbf{SuperPoint + LightGlue}) zur Verbesserung des Feature-Matching-Moduls
    \item Implementierung der gesamten Pipeline in \textbf{C++} mit \textbf{OpenCV}, \textbf{ONNX}-basierter Modellintegration und Bereitstellung als \textbf{REST-API}
  \end{itemize}

  \vspace{0.5em}

  \textbf{Technologien}: Python, C++, OpenCV, BRISK, FLANN, SuperPoint, LightGlue, ONNX, REST-API

  \vspace{0.5em}

  \newpage
  
  % qr- barcode
  \cvsubsubsection{QR- und Barcode-Erkennungssystem}

  \vspace{0.5em}

  \begin{itemize}
    \item Entwicklung einer automatisierten Pipeline zur Erkennung von \textbf{QR- und Barcodes} in industriellen Bilddaten
    \item Implementierung einer \textbf{Multi-Scale-Erkennung} mit Bildvorverarbeitung und robuster Dekodierung mittels \textbf{OpenCV} und \textbf{ZBar}
    \item Deployment des Erkennungssystems als eigenständige Windows-Anwendung
  \end{itemize}

  \vspace{0.5em}

  \textbf{Technologien}: Python, OpenCV, ZBar, uv, Git

  \vspace{0.5em}
  
  % Mesh-Analyse 
  \cvsubsubsection{Mesh-Analyse für Bauwerksmodelle}

  \vspace{0.5em}

  \begin{itemize}
    \item Entwicklung einer Pipeline zur \textbf{2D-Projektion} von \textbf{IFC-Bauwerksmodellen} einschließlich Datenextraktion, Koordinatentransformation und höhenbasierter Filterung
    \item Extraktion und Projektion von \textbf{IfcSpace}-, \textbf{IfcWall}-, \textbf{IfcDoor}- und \textbf{IfcWindow}-Objekten sowie Ausgabe der Polygon- und Objektinformationen im \textbf{JSON}-Format
    \item Implementierung einer \textbf{C\#}-Konsolenanwendung zur internen Verarbeitung und Analyse von Bauwerksmodellen
  \end{itemize}

  \vspace{0.5em}

  \textbf{Technologien}: Python, C\#, IfcOpenShell, Open3D, uv

  \vspace{0.5em}

  % Research & Technology Exploration
  \cvsubsubsection{Forschungs- und Technologieprojekte}

  \vspace{0.5em}

  \begin{itemize}
    \item \textbf{3D-Rekonstruktion — Technologieevaluierung}
          \begin{itemize}
            \item Evaluierung von \textbf{NeRF} und \textbf{3D Gaussian Splatting} für die Rekonstruktion industrieller Szenen, einschließlich lokaler Bereitstellung, Einrichtung der Entwicklungsumgebung sowie Validierung auf internen Datensätzen für Kundendemonstrationen
          \end{itemize}
    \item \textbf{Graphbasierte Dateibeziehungsanalyse}
          \begin{itemize}
            \item Entwicklung eines graphbasierten Prototyps zur Navigation in großen Projektstrukturen in \textbf{C\#}, einschließlich \textbf{Graphaufbau} und \textbf{Suchalgorithmen} zur effizienten Analyse von Dateibeziehungen
          \end{itemize}
    \item \textbf{Interaktive Canvas-Rendering-Engine}
          \begin{itemize}
            \item Erweiterung einer \textbf{TypeScript}-basierten \textbf{Canvas-Rendering-Engine} innerhalb einer \textbf{Blazor-Server}-Anwendung durch Integration neuer Visualisierungskomponenten in das bestehende Rendering-Framework unter Beibehaltung der Kompatibilität mit vorhandenen Funktionen
          \end{itemize}
  \end{itemize}


  \cvsection{AUSBILDUNG}

  \cvevent{M.Sc. Informatik im Bauwesen}{Technische Universität Darmstadt}{04/2018 -- 03/2021}{Darmstadt}

  \cvsubsection{Schwerpunkte}
  \begin{itemize}
    \item Revit-Plugin-Entwicklung mit \textbf{.NET Framework}
    \item BIM-Datenbanken (CRUD) mit \textbf{SQLite}
    \item BIM-Projektmanagement mit \textbf{MS Project}
  \end{itemize}

  \vspace{0.5em}

  \cvsubsection{Wissenschaftliche Hilfskraft (HiWi)}

  SCOPE (Semantic Construction Project Engineering)
  \begin{itemize}
    \item Entwicklung eines webbasierten ViewCube mit Three.js
    \item Ontologiemodellierung der IFC-Hierarchie mit Protégé
  \end{itemize}

  \divider

  \cvevent{B.Sc. Brückenbau}{Chan'an Universität}{09/2012 -- 06/2016}{China}

  % \divider

  \medskip


\end{paracol}


\end{document}
